\documentclass[11pt,a4paper]{article}
\usepackage{amsmath,amssymb,amsthm}
\usepackage[margin=2.4cm]{geometry}
\newtheorem*{theorem}{Theorem}
\newtheorem{lemma}{Lemma}
\newcommand{\abs}[1]{\lvert #1\rvert}
\emergencystretch=1.5em
\title{\large\bfseries A bound for $t$-$(v,k,1)$ packings with $t\ge 3$}
\author{Jiang Chang}
\date{November 2026}
\begin{document}
\maketitle

\paragraph{Notation.}
Throughout, $v\ge k\ge t\ge 2$ are integers, and $V$ is a set of $v$ points.
Put $Q=\binom{k}{t}$ and $b=\binom{k-1}{t-1}$, so that $kb=tQ$.

A \emph{$t$-$(v,k,1)$ packing} is a collection $P$ of $k$-subsets of $V$
(\emph{blocks}) such that every $t$-subset of $V$ lies in at most one block.
Two blocks then meet in at most $t-1$ points, that is, they are at Hamming
distance at least $2(k-t+1)$. So the maximum of $\abs{P}$ is
$A(v,2(k-t+1),k)$, where $A(n,d,w)$ is the maximum size of a binary code of
length $n$, constant weight $w$ and minimum distance at least $d$.

The \emph{leave} of $P$ is the $t$-uniform hypergraph $L$ on $V$ whose edges
are the $t$-subsets of $V$ that lie in no block.

\begin{theorem}
Let $t\ge 3$. Suppose that $\binom{k-s}{t-s}$ divides $\binom{v-s}{t-s}$ for
$s=1,\dots,t-1$, and that $\binom{k}{t}$ does not divide $\binom{v}{t}$. Then
\[
\abs{P}\;\le\;\left\lfloor \binom{v}{t}\Big/\binom{k}{t}\right\rfloor-2 .
\]
\end{theorem}

\paragraph{Examples.}
Let $t=4$ and $k=7$, so that $\binom63=20$, $\binom52=10$, $\binom41=4$ and
$Q=\binom74=35$.

For $v=27$: $\binom{26}{3}=2600=130\cdot 20$, $\binom{25}{2}=300=30\cdot 10$,
$\binom{24}{1}=24=6\cdot 4$, and $\binom{27}{4}=17550=501\cdot 35+15$.
Hence $A(27,8,7)\le 499$.

For $v=47$: $\binom{46}{3}=15180=759\cdot 20$, $\binom{45}{2}=990=99\cdot 10$,
$\binom{44}{1}=44=11\cdot 4$, and $\binom{47}{4}=178365=5096\cdot 35+5$.
Hence $A(47,8,7)\le 5094$.

\medskip
The reason is this. In the leave, the number of edges on a set of $s$
points, $1\le s\le t-1$, is divisible as it is in a complete hypergraph on
$k$ points. A hypergraph with this property is empty, or complete on $k$
points, or has at least $2Q$ edges.

\paragraph{1. The degrees of the leave.}
For a set $S$ of points, $\deg(S)$ is the number of edges containing $S$.
Let $1\le\abs{S}=s\le t-1$. There are $\binom{v-s}{t-s}$ $t$-subsets of $V$
that contain $S$. A block containing $S$ contains $\binom{k-s}{t-s}$ of them,
and no $t$-subset lies in two blocks. So in the leave
\[
\deg(S)=\binom{v-s}{t-s}-m\binom{k-s}{t-s},
\]
where $m$ is the number of blocks containing $S$. By assumption this is
divisible by $\binom{k-s}{t-s}$. So the leave has the property
\begin{itemize}
\item[(D)] $\deg(S)$ is divisible by $\binom{k-s}{t-s}$ whenever
$1\le\abs{S}=s\le t-1$.
\end{itemize}
Each block contains $Q$ $t$-subsets, and no $t$-subset lies in two blocks.
So the leave has $e=\binom{v}{t}-\abs{P}\,Q$ edges.

\paragraph{2. Hypergraphs with property (D).}

\begin{lemma}
Let $k\ge t\ge 2$, and let $L$ be a $t$-uniform hypergraph without repeated
edges, with $e$ edges, satisfying {\rm (D)}. Then every point lies on at most
$te/k$ edges. If $e>0$ then $e\ge Q$, and if $e=Q$ then $L$ is the complete
$t$-uniform hypergraph on $k$ points (and isolated points).
\end{lemma}

\begin{proof}
Write $d(x)=\deg(\{x\})$. By (D) it is a multiple of $b$. The sum of $d(x)$
over all points is $te$. We show $d(x)\le te/k$ by induction on $t$.

\emph{The case $t=2$.} $L$ is a graph in which every degree is a multiple of
$k-1$. Let $x$ have degree $D>0$. Each neighbour of $x$ has positive degree,
hence degree at least $k-1$. The degrees of $x$ and of its $D$ neighbours add
up to at least $D+D(k-1)=kD$, and to at most $2e$. So $D\le 2e/k$.

\emph{The case $t\ge 3$.} For a point $u$, let $L_u$ be the $(t-1)$-uniform
hypergraph whose edges are the sets $E\setminus\{u\}$, for the edges $E$ of
$L$ on $u$. It has $d(u)$ edges, none repeated. If $S$ does not contain $u$
and $1\le\abs{S}=s\le t-2$, then $S$ lies on $\deg(S\cup\{u\})$ edges of
$L_u$, and by (D) this number is divisible by $\binom{k-s-1}{t-s-1}$. So
$L_u$ satisfies (D) with $t-1$, $k-1$ instead of $t$, $k$. By induction, a
point $y$ other than $u$ lies on at most $(t-1)\,d(u)/(k-1)$ edges of $L_u$:
\begin{equation}
\deg(\{u,y\})\le \frac{(t-1)\,d(u)}{k-1}. \tag{1}
\end{equation}
Fix $y$ and add (1) over all $u$ other than $y$. On the left, each edge on
$y$ is counted $t-1$ times, so the sum is $(t-1)\,d(y)$. On the right, the
sum of $d(u)$ over $u$ other than $y$ is $te-d(y)$. Hence
$(k-1)\,d(y)\le te-d(y)$, that is, $d(y)\le te/k$.

\emph{The last two statements.} If $e>0$, some point has $d(x)>0$, hence
$d(x)\ge b$, and $e\ge kb/t=Q$. If $e=Q$, then $d(x)\le tQ/k=b$ for all $x$,
so every point on an edge has $d(x)=b$, and there are $te/b=k$ such points.
The $Q$ edges are distinct $t$-subsets of these $k$ points, and there are
only $\binom{k}{t}=Q$ such subsets. So every one is an edge.
\end{proof}

\begin{lemma}
Let $t\ge 3$, and let $L$ be as in Lemma~1. Then $e=0$, $e=Q$, or $e\ge 2Q$.
\end{lemma}

\begin{proof}
Suppose $0<e<2Q$. For a point $x$ on an edge, let $N(x)$ be the set of the
other points that lie on an edge with $x$.

(a) Let $x$ be on an edge. By Lemma~1, $d(x)\le te/k<2b$, so $d(x)=b$. Then
$L_x$ has $b=\binom{k-1}{t-1}$ edges and satisfies (D) with $t-1\ge 2$ and
$k-1$. By the last part of Lemma~1, $L_x$ is the complete $(t-1)$-uniform
hypergraph on $k-1$ points. These points are $N(x)$. So $\abs{N(x)}=k-1$,
and every $(t-1)$-subset of $N(x)$, with $x$ added, is an edge.

(b) Let $x$ be on an edge and $W=N(x)\cup\{x\}$. Then
$N(y)=W\setminus\{y\}$ for every $y\in W$. For $y=x$ this is the definition
of $W$. Let $y\in N(x)$. Every other point $z$ of $W$ is on an edge with
$y$: if $z=x$, by the definition of $N(x)$; if $z\in N(x)$, take a
$(t-1)$-subset of $N(x)$ containing $y$ and $z$ (possible, as
$2\le t-1\le k-1$) and add $x$. So $N(y)$ contains $W\setminus\{y\}$. By (a)
both have $k-1$ points, so they are equal.

(c) By (b), an edge on a point $y$ of $W$ lies inside $W$, since its other
points are in $N(y)$. And every $t$-subset $T$ of $W$ is an edge: take
$y\in T$; then $T\setminus\{y\}$ is a $(t-1)$-subset of $N(y)$, and by (a)
$T$ is an edge. So $W$ carries a complete hypergraph with $Q$ edges, and no
other edge meets $W$.

(d) If some point $x'$ on an edge is outside $W$, form
$W'=N(x')\cup\{x'\}$ in the same way. If $W'$ met $W$, the complete
hypergraph on $W'$ would have an edge containing $x'$ and a point of $W$,
against (c). So $W'$ is disjoint from $W$. Continuing, the edges of $L$ fall
into complete hypergraphs on disjoint $k$-sets, each with $Q$ edges. Since
$0<e<2Q$ there is exactly one, and $e=Q$.
\end{proof}

\begin{proof}[Proof of the Theorem]
By Section~1, the leave satisfies (D) and has $e=\binom{v}{t}-\abs{P}\,Q$
edges. Let $r$ be the remainder of $\binom{v}{t}$ on division by $Q$; by
assumption $1\le r\le Q-1$. Then $e=aQ+r$ for an integer $a$, so $e$ is
neither $0$ nor $Q$. By Lemma~2, $e\ge 2Q$, hence $a\ge 2$ and
$e\ge 2Q+r$. Therefore
\[
\abs{P}=\frac{\binom{v}{t}-e}{Q}\;\le\;\frac{\binom{v}{t}-r}{Q}-2
=\left\lfloor \binom{v}{t}\Big/Q\right\rfloor-2 . \qedhere
\]
\end{proof}

\paragraph{Remarks.}
Lemma~2 fails for $t=2$: the $4$-cycle has all degrees divisible by $2$
(here $k=3$) and has $4$ edges, between $Q=3$ and $2Q=6$.

For $t=3$, Javadi, Poorhadi and Fallah~\cite{JPF} prove the bound with $-3$
instead of $-2$ under the same conditions.

The argument is close to pages 89--91 of~\cite{B79}.

\begin{thebibliography}{9}\small
\bibitem{B79} A. E. Brouwer, Packing and covering of $\binom{k}{t}$-sets,
in: A. Schrijver (ed.), \emph{Packing and Covering in Combinatorics},
Math. Centre Tracts 106, Amsterdam, 1979, pp.~89--97.
\bibitem{JPF} R. Javadi, E. Poorhadi and F. Fallah, Packing cliques in
$3$-uniform hypergraphs, \emph{J. Combin. Des.} \textbf{28} (2020) 580--603;
Theorem~13 in arXiv:1905.10807.
\end{thebibliography}
\end{document}
